\IfFileExists{stacks-project.cls}{%
\documentclass{stacks-project}
}{%
\documentclass{amsart}
}

% For dealing with references we use the comment environment
\usepackage{verbatim}
\newenvironment{reference}{\comment}{\endcomment}
%\newenvironment{reference}{}{}
\newenvironment{slogan}{\comment}{\endcomment}
\newenvironment{history}{\comment}{\endcomment}

% For commutative diagrams we use Xy-pic
\usepackage[all]{xy}

% We use 2cell for 2-commutative diagrams.
\xyoption{2cell}
\UseAllTwocells

% We use multicol for the list of chapters between chapters
\usepackage{multicol}

% This is generally recommended for better output
\usepackage{lmodern}
\usepackage[T1]{fontenc}

% For cross-file-references
\usepackage{xr-hyper}

% Package for hypertext links:
\usepackage{hyperref}

% For any local file, say "hello.tex" you want to link to please
% use \externaldocument[hello-]{hello}
\externaldocument[introduction-]{introduction}
\externaldocument[conventions-]{conventions}
\externaldocument[sets-]{sets}
\externaldocument[categories-]{categories}
\externaldocument[topology-]{topology}
\externaldocument[sheaves-]{sheaves}
\externaldocument[sites-]{sites}
\externaldocument[stacks-]{stacks}
\externaldocument[fields-]{fields}
\externaldocument[algebra-]{algebra}
\externaldocument[brauer-]{brauer}
\externaldocument[homology-]{homology}
\externaldocument[derived-]{derived}
\externaldocument[simplicial-]{simplicial}
\externaldocument[more-algebra-]{more-algebra}
\externaldocument[smoothing-]{smoothing}
\externaldocument[modules-]{modules}
\externaldocument[sites-modules-]{sites-modules}
\externaldocument[injectives-]{injectives}
\externaldocument[cohomology-]{cohomology}
\externaldocument[sites-cohomology-]{sites-cohomology}
\externaldocument[dga-]{dga}
\externaldocument[dpa-]{dpa}
\externaldocument[sdga-]{sdga}
\externaldocument[hypercovering-]{hypercovering}
\externaldocument[schemes-]{schemes}
\externaldocument[constructions-]{constructions}
\externaldocument[properties-]{properties}
\externaldocument[morphisms-]{morphisms}
\externaldocument[coherent-]{coherent}
\externaldocument[divisors-]{divisors}
\externaldocument[limits-]{limits}
\externaldocument[varieties-]{varieties}
\externaldocument[topologies-]{topologies}
\externaldocument[descent-]{descent}
\externaldocument[perfect-]{perfect}
\externaldocument[more-morphisms-]{more-morphisms}
\externaldocument[flat-]{flat}
\externaldocument[groupoids-]{groupoids}
\externaldocument[more-groupoids-]{more-groupoids}
\externaldocument[etale-]{etale}
\externaldocument[chow-]{chow}
\externaldocument[intersection-]{intersection}
\externaldocument[pic-]{pic}
\externaldocument[weil-]{weil}
\externaldocument[adequate-]{adequate}
\externaldocument[dualizing-]{dualizing}
\externaldocument[duality-]{duality}
\externaldocument[discriminant-]{discriminant}
\externaldocument[derham-]{derham}
\externaldocument[local-cohomology-]{local-cohomology}
\externaldocument[algebraization-]{algebraization}
\externaldocument[curves-]{curves}
\externaldocument[resolve-]{resolve}
\externaldocument[models-]{models}
\externaldocument[functors-]{functors}
\externaldocument[equiv-]{equiv}
\externaldocument[pione-]{pione}
\externaldocument[etale-cohomology-]{etale-cohomology}
\externaldocument[proetale-]{proetale}
\externaldocument[relative-cycles-]{relative-cycles}
\externaldocument[more-etale-]{more-etale}
\externaldocument[trace-]{trace}
\externaldocument[crystalline-]{crystalline}
\externaldocument[spaces-]{spaces}
\externaldocument[spaces-properties-]{spaces-properties}
\externaldocument[spaces-morphisms-]{spaces-morphisms}
\externaldocument[decent-spaces-]{decent-spaces}
\externaldocument[spaces-cohomology-]{spaces-cohomology}
\externaldocument[spaces-limits-]{spaces-limits}
\externaldocument[spaces-divisors-]{spaces-divisors}
\externaldocument[spaces-over-fields-]{spaces-over-fields}
\externaldocument[spaces-topologies-]{spaces-topologies}
\externaldocument[spaces-descent-]{spaces-descent}
\externaldocument[spaces-perfect-]{spaces-perfect}
\externaldocument[spaces-more-morphisms-]{spaces-more-morphisms}
\externaldocument[spaces-flat-]{spaces-flat}
\externaldocument[spaces-groupoids-]{spaces-groupoids}
\externaldocument[spaces-more-groupoids-]{spaces-more-groupoids}
\externaldocument[bootstrap-]{bootstrap}
\externaldocument[spaces-pushouts-]{spaces-pushouts}
\externaldocument[spaces-chow-]{spaces-chow}
\externaldocument[groupoids-quotients-]{groupoids-quotients}
\externaldocument[spaces-more-cohomology-]{spaces-more-cohomology}
\externaldocument[spaces-simplicial-]{spaces-simplicial}
\externaldocument[spaces-duality-]{spaces-duality}
\externaldocument[formal-spaces-]{formal-spaces}
\externaldocument[restricted-]{restricted}
\externaldocument[spaces-resolve-]{spaces-resolve}
\externaldocument[formal-defos-]{formal-defos}
\externaldocument[defos-]{defos}
\externaldocument[cotangent-]{cotangent}
\externaldocument[examples-defos-]{examples-defos}
\externaldocument[algebraic-]{algebraic}
\externaldocument[examples-stacks-]{examples-stacks}
\externaldocument[stacks-sheaves-]{stacks-sheaves}
\externaldocument[criteria-]{criteria}
\externaldocument[artin-]{artin}
\externaldocument[quot-]{quot}
\externaldocument[stacks-properties-]{stacks-properties}
\externaldocument[stacks-morphisms-]{stacks-morphisms}
\externaldocument[stacks-limits-]{stacks-limits}
\externaldocument[stacks-cohomology-]{stacks-cohomology}
\externaldocument[stacks-perfect-]{stacks-perfect}
\externaldocument[stacks-introduction-]{stacks-introduction}
\externaldocument[stacks-more-morphisms-]{stacks-more-morphisms}
\externaldocument[stacks-geometry-]{stacks-geometry}
\externaldocument[moduli-]{moduli}
\externaldocument[moduli-curves-]{moduli-curves}
\externaldocument[examples-]{examples}
\externaldocument[exercises-]{exercises}
\externaldocument[guide-]{guide}
\externaldocument[desirables-]{desirables}
\externaldocument[coding-]{coding}
\externaldocument[obsolete-]{obsolete}
\externaldocument[fdl-]{fdl}
\externaldocument[index-]{index}

% Theorem environments.
%
\theoremstyle{plain}
\newtheorem{theorem}[subsection]{Theorem}
\newtheorem{proposition}[subsection]{Proposition}
\newtheorem{lemma}[subsection]{Lemma}

\theoremstyle{definition}
\newtheorem{definition}[subsection]{Definition}
\newtheorem{example}[subsection]{Example}
\newtheorem{exercise}[subsection]{Exercise}
\newtheorem{situation}[subsection]{Situation}

\theoremstyle{remark}
\newtheorem{remark}[subsection]{Remark}
\newtheorem{remarks}[subsection]{Remarks}

\numberwithin{equation}{subsection}

% Macros
%
\def\lim{\mathop{\mathrm{lim}}\nolimits}
\def\colim{\mathop{\mathrm{colim}}\nolimits}
\def\Spec{\mathop{\mathrm{Spec}}}
\def\Hom{\mathop{\mathrm{Hom}}\nolimits}
\def\Ext{\mathop{\mathrm{Ext}}\nolimits}
\def\SheafHom{\mathop{\mathcal{H}\!\mathit{om}}\nolimits}
\def\SheafExt{\mathop{\mathcal{E}\!\mathit{xt}}\nolimits}
\def\Sch{\mathit{Sch}}
\def\Mor{\mathop{\mathrm{Mor}}\nolimits}
\def\Ob{\mathop{\mathrm{Ob}}\nolimits}
\def\Sh{\mathop{\mathit{Sh}}\nolimits}
\def\NL{\mathop{N\!L}\nolimits}
\def\CH{\mathop{\mathrm{CH}}\nolimits}
\def\proetale{{pro\text{-}\acute{e}tale}}
\def\etale{{\acute{e}tale}}
\def\QCoh{\mathit{QCoh}}
\def\Ker{\mathop{\mathrm{Ker}}}
\def\Im{\mathop{\mathrm{Im}}}
\def\Coker{\mathop{\mathrm{Coker}}}
\def\Coim{\mathop{\mathrm{Coim}}}

% Boxtimes
%
\DeclareMathSymbol{\boxtimes}{\mathbin}{AMSa}{"02}

%
% Macros for moduli stacks/spaces
%
\def\QCohstack{\mathcal{QC}\!\mathit{oh}}
\def\Cohstack{\mathcal{C}\!\mathit{oh}}
\def\Spacesstack{\mathcal{S}\!\mathit{paces}}
\def\Quotfunctor{\mathrm{Quot}}
\def\Hilbfunctor{\mathrm{Hilb}}
\def\Curvesstack{\mathcal{C}\!\mathit{urves}}
\def\Polarizedstack{\mathcal{P}\!\mathit{olarized}}
\def\Complexesstack{\mathcal{C}\!\mathit{omplexes}}
% \Pic is the operator that assigns to X its picard group, usage \Pic(X)
% \Picardstack_{X/B} denotes the Picard stack of X over B
% \Picardfunctor_{X/B} denotes the Picard functor of X over B
\def\Pic{\mathop{\mathrm{Pic}}\nolimits}
\def\Picardstack{\mathcal{P}\!\mathit{ic}}
\def\Picardfunctor{\mathrm{Pic}}
\def\Deformationcategory{\mathcal{D}\!\mathit{ef}}
\usepackage{cleveref}
\crefname{originalreference}{original reference}{original references}
\Crefname{originalreference}{Original reference}{Original references}
\makeatletter
\newlabel{categories-definition-cofinal}{{04E6}{source}{}{sourcepointer-1}{}}
\newlabel{categories-definition-cofinal@cref}{{[originalreference][0][]04E6}{[1][0][]source}{}{}{}}
\makeatother
\crefname{originalreference}{original reference}{original references}
\Crefname{originalreference}{Original reference}{Original references}
\makeatletter
\newlabel{categories-lemma-cofinal-essentially-constant}{{0A1S}{source}{}{sourcepointer-2}{}}
\newlabel{categories-lemma-cofinal-essentially-constant@cref}{{[originalreference][0][]0A1S}{[1][0][]source}{}{}{}}
\makeatother
\crefname{originalreference}{original reference}{original references}
\Crefname{originalreference}{Original reference}{Original references}
\makeatletter
\newlabel{categories-lemma-essentially-constant-over-product}{{0A2C}{source}{}{sourcepointer-3}{}}
\newlabel{categories-lemma-essentially-constant-over-product@cref}{{[originalreference][0][]0A2C}{[1][0][]source}{}{}{}}
\makeatother
\crefname{originalreference}{original reference}{original references}
\Crefname{originalreference}{Original reference}{Original references}
\makeatletter
\newlabel{categories-lemma-what-gets-inverted}{{05Q9}{source}{}{sourcepointer-4}{}}
\newlabel{categories-lemma-what-gets-inverted@cref}{{[originalreference][0][]05Q9}{[1][0][]source}{}{}{}}
\makeatother
\crefname{originalreference}{original reference}{original references}
\Crefname{originalreference}{Original reference}{Original references}
\makeatletter
\newlabel{categories-remark-left-localization-morphisms-colimit}{{05Q0}{source}{}{sourcepointer-5}{}}
\newlabel{categories-remark-left-localization-morphisms-colimit@cref}{{[originalreference][0][]05Q0}{[1][0][]source}{}{}{}}
\makeatother
\crefname{originalreference}{original reference}{original references}
\Crefname{originalreference}{Original reference}{Original references}
\makeatletter
\newlabel{homology-lemma-direct-sum-from-product-essentially-constant}{{0A2G}{source}{}{sourcepointer-6}{}}
\newlabel{homology-lemma-direct-sum-from-product-essentially-constant@cref}{{[originalreference][0][]0A2G}{[1][0][]source}{}{}{}}
\makeatother
\crefname{originallemma}{original lemma}{original lemmas}
\Crefname{originallemma}{Original lemma}{Original lemmas}
\makeatletter
\newlabel{lemma-2-out-of-3-defined}{{05SC}{source}{}{sourcepointer-7}{}}
\newlabel{lemma-2-out-of-3-defined@cref}{{[originallemma][0][]05SC}{[1][0][]source}{}{}{}}
\makeatother
\crefname{originallemma}{original lemma}{original lemmas}
\Crefname{originallemma}{Original lemma}{Original lemmas}
\makeatletter
\newlabel{lemma-split}{{05QT}{source}{}{sourcepointer-8}{}}
\newlabel{lemma-split@cref}{{[originallemma][0][]05QT}{[1][0][]source}{}{}{}}
\makeatother
\crefname{originalreference}{original reference}{original references}
\Crefname{originalreference}{Original reference}{Original references}
\makeatletter
\newlabel{situation-derived-functor}{{05S8}{source}{}{sourcepointer-9}{}}
\newlabel{situation-derived-functor@cref}{{[originalreference][0][]05S8}{[1][0][]source}{}{}{}}
\makeatother
\title{RF is defined at a direct sum exactly when defined at both summands}
\author{The Stacks Project Authors}
\date{stacks-project:derived.tex:local-snapshot-20261012; GNU FDL 1.2 or later}
\begin{document}
\maketitle
\noindent\textit{Editorial note.} Counted span is the author's own proof environment, source lines 5031-5109, following the statement at 5018-5029.\par\medskip
\noindent\textit{Editorial note.} The closing 'The proof of the corresponding statements for LF are dual' (line 5108) is a mirror-level omission, recorded.\par\medskip
\noindent\textit{Allowed standard prerequisites (editorial).}\begin{itemize}
\item Split triangles. lemma-split, cited at proof lines 5033 and 5038; complete proof at source lines 916-927, queued for separate carding. Supplies the split triangle X -> X + Y -> Y -> X[1] and the identification of the middle triangle with a direct sum.
\item Two-out-of-three for definedness. lemma-2-out-of-3-defined (carded in this batch, L26162), cited at proof line 5034. Supplies that RF defined at X and Y implies defined at X + Y with the split triangle.
\item Arrows inverted by localization. categories-lemma-what-gets-inverted, cross-chapter citation at proof line 5047 (tags mechanism). Supplies that saturated S is exactly the set of arrows inverted by Q, giving s + s' in S.
\item Filteredness of slice categories and products. categories-remark-left-localization-morphisms-colimit and categories-lemma-essentially-constant-over-product, cross-chapter citations at proof lines 5055-5058 (tags mechanism). Supplies X/S, Y/S, (X + Y)/S filtered and the product criterion for essential constancy.
\item Direct sums of essentially constant systems. homology-lemma-direct-sum-from-product-essentially-constant, cross-chapter citation at proof line 5070 (tags mechanism; Karoubian hypothesis enters here). Supplies that essential constancy of the direct sum system implies essential constancy of both summand systems.
\end{itemize}\medskip
\ifdefined\hypertarget\hypertarget{sourcepointer-1}{}\fi
\noindent\textit{Source reference.} Original reference 04E6: focused-seeds/geometry/S9-der-direct-sum-defined/ref-tags.txt, source lines 5--5. Navigation pointer only: the counted proof cites this statement by \textbackslash{}ref as a standing input; the deduction the counted proof has to make is carried out inside the counted span, and no part of that deduction is delegated to this pointer.\par\medskip
\ifdefined\hypertarget\hypertarget{sourcepointer-2}{}\fi
\noindent\textit{Source reference.} Original reference 0A1S: focused-seeds/geometry/S9-der-direct-sum-defined/ref-tags.txt, source lines 6--6. Navigation pointer only: the counted proof cites this statement by \textbackslash{}ref as a standing input; the deduction the counted proof has to make is carried out inside the counted span, and no part of that deduction is delegated to this pointer.\par\medskip
\ifdefined\hypertarget\hypertarget{sourcepointer-3}{}\fi
\noindent\textit{Source reference.} Original reference 0A2C: focused-seeds/geometry/S9-der-direct-sum-defined/ref-tags.txt, source lines 7--7. Navigation pointer only: the counted proof cites this statement by \textbackslash{}ref as a standing input; the deduction the counted proof has to make is carried out inside the counted span, and no part of that deduction is delegated to this pointer.\par\medskip
\ifdefined\hypertarget\hypertarget{sourcepointer-4}{}\fi
\noindent\textit{Source reference.} Original reference 05Q9: focused-seeds/geometry/S9-der-direct-sum-defined/ref-tags.txt, source lines 8--8. Navigation pointer only: the counted proof cites this statement by \textbackslash{}ref as a standing input; the deduction the counted proof has to make is carried out inside the counted span, and no part of that deduction is delegated to this pointer.\par\medskip
\ifdefined\hypertarget\hypertarget{sourcepointer-5}{}\fi
\noindent\textit{Source reference.} Original reference 05Q0: focused-seeds/geometry/S9-der-direct-sum-defined/ref-tags.txt, source lines 9--9. Navigation pointer only: the counted proof cites this statement by \textbackslash{}ref as a standing input; the deduction the counted proof has to make is carried out inside the counted span, and no part of that deduction is delegated to this pointer.\par\medskip
\ifdefined\hypertarget\hypertarget{sourcepointer-6}{}\fi
\noindent\textit{Source reference.} Original reference 0A2G: focused-seeds/geometry/S9-der-direct-sum-defined/ref-tags.txt, source lines 10--10. Navigation pointer only: the counted proof cites this statement by \textbackslash{}ref as a standing input; the deduction the counted proof has to make is carried out inside the counted span, and no part of that deduction is delegated to this pointer.\par\medskip
\ifdefined\hypertarget\hypertarget{sourcepointer-7}{}\fi
\noindent\textit{Source reference.} Original lemma 05SC: sources/stacks-derived/pinned-original-source.tex, source lines 4896--4896. Navigation pointer only: the counted proof cites this statement by \textbackslash{}ref as a standing input; the deduction the counted proof has to make is carried out inside the counted span, and no part of that deduction is delegated to this pointer.\par\medskip
\ifdefined\hypertarget\hypertarget{sourcepointer-8}{}\fi
\noindent\textit{Source reference.} Original lemma 05QT: sources/stacks-derived/pinned-original-source.tex, source lines 917--917. Navigation pointer only: the counted proof cites this statement by \textbackslash{}ref as a standing input; the deduction the counted proof has to make is carried out inside the counted span, and no part of that deduction is delegated to this pointer.\par\medskip
\ifdefined\hypertarget\hypertarget{sourcepointer-9}{}\fi
\noindent\textit{Source reference.} Original reference 05S8: sources/stacks-derived/pinned-original-source.tex, source lines 4736--4736. Navigation pointer only: the counted proof cites this statement by \textbackslash{}ref as a standing input; the deduction the counted proof has to make is carried out inside the counted span, and no part of that deduction is delegated to this pointer.\par\medskip
\begin{lemma}
\label{lemma-direct-sum-defined}
Assumptions and notation as in Situation \ref{situation-derived-functor}.
Let $X, Y$ be objects of $\mathcal{D}$.
\begin{enumerate}
\item If $RF$ is defined at $X$ and $Y$, then $RF$ is defined at $X \oplus Y$.
\item If $\mathcal{D}'$ is Karoubian and $RF$ is defined at $X \oplus Y$,
then $RF$ is defined at both $X$ and $Y$.
\end{enumerate}
In either case we have $RF(X \oplus Y) = RF(X) \oplus RF(Y)$.
Similarly for $LF$.
\end{lemma}

\begin{proof}
If $RF$ is defined at $X$ and $Y$, then the distinguished triangle
$X \to X \oplus Y \to Y \to X[1]$ (Lemma \ref{lemma-split}) and
Lemma \ref{lemma-2-out-of-3-defined}
shows that $RF$ is defined at $X \oplus Y$ and that we
have a distinguished triangle
$RF(X) \to RF(X \oplus Y) \to RF(Y) \to RF(X)[1]$.
Applying Lemma \ref{lemma-split} to this once more we find
that $RF(X \oplus Y) = RF(X) \oplus RF(Y)$.
This proves (1) and the final assertion.

\medskip\noindent
Conversely, assume that $RF$ is defined at $X \oplus Y$ and that $\mathcal{D}'$
is Karoubian. Since $S$ is a saturated system $S$ is the set of arrows which
become invertible under the additive localization functor
$Q : \mathcal{D} \to S^{-1}\mathcal{D}$, see
Categories, Lemma \ref{categories-lemma-what-gets-inverted}.
Thus for any $s : X \to X'$ and $s' : Y \to Y'$ in $S$ the morphism
$s \oplus s' : X \oplus Y \to X' \oplus Y'$ is an element of $S$.
In this way we obtain a functor
$$
X/S \times Y/S \longrightarrow (X \oplus Y)/S
$$
Recall that the categories $X/S, Y/S, (X \oplus Y)/S$ are filtered
(Categories, Remark
\ref{categories-remark-left-localization-morphisms-colimit}).
By Categories, Lemma \ref{categories-lemma-essentially-constant-over-product}
$X/S \times Y/S$ is filtered and
$F|_{X/S} : X/S \to \mathcal{D}'$ (resp.\ $G|_{Y/S} : Y/S \to \mathcal{D}'$)
is essentially constant if and only if
$F|_{X/S} \circ \text{pr}_1 : X/S \times Y/S \to \mathcal{D}'$
(resp.\ $G|_{Y/S} \circ \text{pr}_2 : X/S \times Y/S \to \mathcal{D}'$)
is essentially constant. Below we will show that the displayed functor
is cofinal, hence by
Categories, Lemma \ref{categories-lemma-cofinal-essentially-constant},
we see that $F|_{(X \oplus Y)/S}$ is essentially constant implies that
$F|_{X/S} \circ \text{pr}_1 \oplus F|_{Y/S} \circ \text{pr}_2 :
X/S \times Y/S \to \mathcal{D}'$
is essentially constant. By Homology, Lemma
\ref{homology-lemma-direct-sum-from-product-essentially-constant}
(and this is where we use that $\mathcal{D}'$ is Karoubian)
we see that 
$F|_{X/S} \circ \text{pr}_1 \oplus F|_{Y/S} \circ \text{pr}_2$
being essentially constant implies
$F|_{X/S} \circ \text{pr}_1$ and
$F|_{Y/S} \circ \text{pr}_2$ are essentially constant proving that $RF$ is
defined at $X$ and $Y$.

\medskip\noindent
Proof that the displayed functor is cofinal.
To do this pick any $t : X \oplus Y \to Z$ in $S$.
Using MS2 we can find morphisms $Z \to X'$, $Z \to Y'$
and $s : X \to X'$, $s' : Y \to Y'$ in $S$ such that
$$
\xymatrix{
X \ar[d]^s & X \oplus Y \ar[d] \ar[l] \ar[r] & Y \ar[d]_{s'} \\
X' & Z \ar[l] \ar[r] & Y'
}
$$
commutes. This proves there is a map $Z \to X' \oplus Y'$ in
$(X \oplus Y)/S$, i.e., we get part (1) of Categories, Definition
\ref{categories-definition-cofinal}. To prove part (2) it suffices
to prove that given $t : X \oplus Y \to Z$ and morphisms
$s_i \oplus s'_i : Z \to X'_i \oplus Y'_i$, $i = 1, 2$ in $(X \oplus Y)/S$
we can find morphisms $a : X'_1 \to X'$, $b : X'_2 \to X'$,
$c : Y'_1 \to Y'$, $d : Y'_2 \to Y'$ in $S$ such that
$a \circ s_1 = b \circ s_2$ and $c \circ s'_1 = d \circ s'_2$.
To do this we first choose any $X'$ and $Y'$ and maps $a, b, c, d$
in $S$; this is possible as $X/S$ and $Y/S$ are filtered. Then the
two maps $a \circ s_1, b \circ s_2 : Z \to X'$ become equal in
$S^{-1}\mathcal{D}$. Hence we can find a morphism
$X' \to X''$ in $S$ equalizing them. Similarly we find $Y' \to Y''$ in $S$
equalizing $c \circ s'_1$ and $d \circ s'_2$. Replacing $X'$ by $X''$ and
$Y'$ by $Y''$ we get $a \circ s_1 = b \circ s_2$ and
$c \circ s'_1 = d \circ s'_2$.

\medskip\noindent
The proof of the corresponding statements for $LF$ are dual.
\end{proof}

\end{document}
